\subsection{Retorno longitudinal y determinación de la escala areal}
\label{rev5:longitud}

La unidad dimensional del área admite una construcción explícita
desde el mismo estado enriquecido. Escribiremos \(L_\alpha\) para la
longitud física obtenida; el entero \(L\) que designa la longitud de un
registro en la sección~\ref{sec:compat-lectores-formas} conserva su
significado combinatorio. La distinción separa el cardinal de una
representación, la extensión de una arista y el radio de su realización.

La composición utiliza las coordenadas regionales y el registro ya
generados por APP--TRIT--TPK. APP conserva ambas hojas con sus residuos
y cocientes; TRIT orienta el régimen local; TPK prolonga el estado,
conservando acarreo, frontera, ruta y memoria. La cadena
\(w_6\to w_{12}\to w_{18}\to w_{24}\to w_{30}\to R_{36}\to G_9\)
y el orden \(U_t\to\Gamma_9\to K_{\rm ph}\) mantienen el
retorno de fase con incremento de memoria. La construcción conjunta
del continuo conserva sus aspectos espectral, solenoidal, gauge,
cohomológico y determinantal. La longitud que sigue es una lectura
dimensional posterior de esa estructura y de sus salidas correlacionadas
\(\pi_{\rm HMT},\varphi_{\rm HMT},e_{\rm HMT},\alpha_{\rm HMT}\).

Fijemos el registro marcado
\[
 \Omega=\{(j,k,q,u,v):j,k\in\mathbb Z/12\mathbb Z,\quad
 k-j\in\{0,3,6,9\},\quad1\le q\le8,\quad
 u<v,\quad u,v\in\{1,2,4,5,7,8\}\}.
\]
La resta se calcula en ese grupo cíclico. Su cardinal es
\(N_\Omega=12\cdot4\cdot8\binom62=5760\).
El factor \(120\) procede del espacio octádico de incidencias
punto--pareja ya construido; las doce posiciones y las cuatro
posiciones de cada órbita conservan sus marcas respectivas. Esta
resolución amplía el portador del retorno de Hadamard. En los espacios
hermitianos de partida y llegada, de dimensión \(N_\Omega\), sea
\(B\) el inversor transportado, con
\(B^*B=BB^*=I/4\). La identidad local
\(H_4^*H_4=4I\) produce esa normalización para el inversor;
su extensión ortogonal y el transporte de las marcas la conservan.

En el espacio de llegada, sean
\(u_\Omega=N_\Omega^{-1/2}\mathbf1\),
\(P=u_\Omega u_\Omega^*\), \(C=(I-P)B\) y \(M=PB\).
Se conserva el transporte completo \(C+M=B\). Los proyectores
atómicos \(P_i=e_ie_i^*\), determinados por la inscripción,
definen la expectativa diagonal
\(\Delta_\Omega(T)=\sum_iP_iTP_i\).

\begin{proposition}[Respuesta registrada y complemento conservado]
\label{rev5:retorno}
Entre las aplicaciones lineales hacia el álgebra diagonal que fijan
los \(P_i\) y son bimodulares respecto de esa álgebra,
\(\Delta_\Omega\) es única. Además,
\begin{equation}
 \Delta_\Omega(CC^*)=r_\Omega I,
 \qquad r_\Omega=\frac{N_\Omega-1}{4N_\Omega}
 =\frac{5759}{23040},
 \qquad \Delta_\Omega(MM^*)=\frac{I}{4N_\Omega}.
 \label{rev5:retorno-escalar}
\end{equation}
Las dos respuestas suman \(I/4\).
\end{proposition}
\begin{proof}
Para la unidad matricial \(E_{ij}\), la bimodularidad impone
\(E(E_{ij})=P_iE(E_{ij})P_j\). Como la imagen es diagonal,
este término vale cero si \(i\ne j\); si \(i=j\), fijar
\(P_i\) determina su imagen. Las unidades matriciales constituyen
una base y determinan el mapa. Por otra parte,
\(CC^*=(I-P)/4\), \(MM^*=P/4\), y cada entrada diagonal
de \(P\) es \(1/N_\Omega\). Aplicar la expectativa demuestra
las tres identidades.
\end{proof}

La inscripción isométrica \(We_i=e_i\otimes e_i\) realiza
\(W^*(T\otimes I)W=\Delta_\Omega(T)\). La inscripción
completa conserva la información del registro; la expectativa extrae
su respuesta diagonal. El complemento \(M\) permanece en el
transporte y posee la respuesta específica de
\eqref{rev5:retorno-escalar}. Esta separación identifica el observable
que interviene en la longitud y el componente complementario conservado.

La forma normal de Smith de \(H_4\) es
\(\operatorname{diag}(1,2,2,4)\); su conúcleo tiene orden
dieciséis. El producto del escalar generado
\(\alpha=\alpha_{\rm HMT}\), indexado por las clases de ese
conúcleo, define el lector multiplicativo finito
\[
 n_H:=\prod_{\bar v\in\operatorname{coker}_{\mathbb Z}(H_4)}
 \alpha=\alpha^{16}.
\]
Sea \(L_*>0\) una sección de referencia
radial de la recta dimensional de longitud. La composición longitudinal
actúa linealmente sobre una cocadena orientada \(\beta_*\):
\begin{equation}
 \mathscr T_L\beta_*
 =\alpha^{16}(\Delta_\Omega(CC^*)\otimes I)\beta_*
 =q_\alpha\beta_*,\qquad
 q_\alpha=\alpha^{16}r_\Omega,\qquad
 L_\alpha=q_\alpha L_*.
 \label{rev5:longitud-generada}
\end{equation}
La fórmula fija una composición concreta de lectores. La linealidad
conserva el tipo de longitud y su orientación; el coeficiente procede
del retorno registrado y de la norma local. Por ejemplo, la respuesta
diagonal de una fuente \(a e_i\) es \(a r_\Omega\), mientras
la norma \(\|C^*ae_i\|=|a|\sqrt{r_\Omega}\) constituye
otra lectura del mismo transporte.

Si \(\beta_*(\bar a)=-U(a)\beta_*(a)\), multiplicar por
\(q_\alpha\) conserva la inversión orientada. Si una arista se
subdivide en nueve tramos compatibles,
\(\beta_*(a)=\sum_jU_j^{-1}\beta_*(a_j)\), la misma
multiplicación conserva esa identidad y la concatenación de memoria.
En una ampliación auxiliar del registro, se transportan
\(P\), \(CC^*\) y \(\Delta_\Omega\) por tensorización
con la identidad. Así permanece \(r_\Omega\); el cardinal de
una malla refinada y el cardinal del registro marcado tienen funciones
distintas.

\subsection{Reciprocidad radial y reconocimiento de la longitud de Planck}
\label{rev5:radial}

La lectura circular del calendario de ciento ocho avances conserva
\(108\ell_0=2\pi_{\rm HMT}L_\alpha\), con
\(\ell_0=ct_0\). Por consiguiente,
\begin{equation}
 t_0=\frac{\pi_{\rm HMT}L_\alpha}{54c},\qquad
 \omega=\frac{c}{L_\alpha},\qquad
 Q_{\rm clk}=\hbar_{\rm ret}\omega
 =\frac{\hbar_{\rm ret}c}{L_\alpha}.
 \label{rev5:reloj}
\end{equation}
La acción retornada se recibe de su ecuación HMT completa en
\eqref{eq:mass-k-action}; la velocidad se recibe de los dos canales
constitutivos de la sección~\ref{sec:compat-lectores-formas}.
La relación circular convierte el avance levantado en radio. Si
\(c=\widehat c\,u_C\), \(t_0=u_T\) y
\(u_L=u_Cu_T\), resulta
\(L_*/u_L=54\widehat c/(\pi_{\rm HMT}q_\alpha)\).
La referencia radial \(L_*\) y la base de arista \(u_L\)
quedan relacionadas con sus tipos dimensionales explícitos.

En una fibra finita, sea \(X=X^*>0\) un carácter positivo
invertible, \(U=2B\) y \(Y=UXU^*\). El transporte
longitudinal es \(D_L=L_\alpha U\). Su respuesta radial positiva
se define mediante la métrica de llegada,
\(\|R_Xv\|^2=\|XD_L^*v\|^2\) para todo \(v\).
La energía y la longitud conjugada conservan el mismo carácter:
\[
 H_X=Q_{\rm clk}Y,\qquad M_X=H_X/c^2,\qquad
 \Lambda_X=L_\alpha Y^{-1}.
\]

\begin{theorem}[Coeficiente radial común y escala areal]
\label{rev5:rigidez-radial}
Las condiciones anteriores determinan una única respuesta positiva
\(R_X=L_\alpha Y\). Se cumplen
\begin{equation}
 R_X\Lambda_X=L_\alpha^2I,\qquad
 H_X\Lambda_X=\hbar_{\rm ret}cI,\qquad
 R_X=\frac{L_\alpha^2}{\hbar_{\rm ret}c}H_X.
\end{equation}
En la realización física que identifica \(R_X\) con el radio
gravitatorio reducido \(R_X=(G/c^2)M_X\), el coeficiente es
único y común a todos los caracteres positivos admisibles:
\begin{equation}
 G=\frac{c^3L_\alpha^2}{\hbar_{\rm ret}},\qquad
 \frac{\hbar_{\rm ret}G}{c^3}=L_\alpha^2.
 \label{rev5:G-area}
\end{equation}
\end{theorem}
\begin{proof}
La polarización de la igualdad de normas determina
\(R_X^2=D_LX^2D_L^*=L_\alpha^2UX^2U^*\).
El operador \(L_\alpha UXU^*\) es positivo y su cuadrado es
el indicado. La unicidad de la raíz positiva prueba la primera
afirmación. Las dos identidades de productos siguen por sustitución.
La convención física da
\(L_\alpha Y=(GQ_{\rm clk}/c^4)Y\). La invertibilidad
de \(Y\) permite cancelarlo, y la expresión de
\(Q_{\rm clk}\) proporciona \eqref{rev5:G-area}.
Si otro coeficiente conserva los mismos datos, su diferencia
multiplicada por \(M_X\) vale cero y, por tanto, es nula.
\end{proof}

El dominio de la prueba es la fibra positiva finita; también admite
operadores acotados uniformemente positivos con las mismas identidades.
La lectura de radio gravitatorio reducido constituye la convención
física declarada de esta realización. El reconocimiento convencional
posterior \(\ell_P^2=\hbar_{\rm ret}G/c^3\) identifica
\(\ell_P=L_\alpha\). La longitud se obtiene primero mediante
\eqref{rev5:longitud-generada}; esta identificación conserva su
constructor. La fórmula circular equivalente es
\[
 G=\left(\frac{54}{\pi_{\rm HMT}}\right)^2
 \frac{c^5t_0^2}{\hbar_{\rm ret}}.
\]
El factor circular ya está incluido cuando se utiliza \(L_\alpha\).
La longitud y la sección de acción determinan asimismo
\[
 t_P=\frac{L_\alpha}{c},\qquad
 m_P=\frac{\hbar_{\rm ret}}{cL_\alpha},\qquad
 E_P=\frac{\hbar_{\rm ret}c}{L_\alpha}.
\]
Estas identidades siguen al sustituir \eqref{rev5:G-area} en las
expresiones convencionales de las tres escalas; la positividad
determina sus raíces. Para un modo de carácter \(y>0\), las
lecturas son \(m(y)=m_Py\), \(r_g(y)=L_\alpha y\) y
\(\bar\lambda(y)=L_\alpha/y\). Su producto longitudinal
vale \(r_g(y)\bar\lambda(y)=L_\alpha^2\). La igualdad de
ambas longitudes equivale a \(y=y^{-1}\) y, por positividad,
caracteriza exactamente \(y=1\). El calendario conserva
\(t_0=(\pi_{\rm HMT}/54)t_P\) y
\(108t_0=2\pi_{\rm HMT}t_P\).

Un radio de horizonte \(R_h=L_\alpha\zeta_h\) conserva su
factor propio \(\zeta_h=R_h/L_\alpha\). Su área espacial
pertenece al horizonte declarado; la escala elemental del operador
sectorial siguiente es \(L_\alpha^2\). En la especialización
de Schwarzschild para la misma masa, el radio es \(2r_g(y)\)
y, por tanto, \(\zeta_h=2y\); un horizonte de otra realización
conserva su condición geométrica de formación.
